\documentclass[quiet]{ctexart}

\ExplSyntaxOn \makeatletter

\keys_define:nn { hduthesis / exam }
  {
    watermark .tl_set:N = \l__hduthesis_exam_watermark_tl,
  }
\NewDocumentCommand \hdutexset {m} { \keys_set:nn { hduthesis / exam } {#1} }
\DeclareHookRule { shipout / background }
  { hduthesis / exam } { before } { pgfrcs }
\AddToHook { shipout / background } [ hduthesis / exam ]
  {
    \tl_if_empty:NF \l__hduthesis_exam_watermark_tl
      {
        \tikz [ remember~picture, overlay ]
          \node [ opacity = .2 ] at (current~page)
            {
              \includegraphics [ width = \linewidth ]
                { \l__hduthesis_exam_watermark_tl }
            };
      }
  }

\RequirePackage
  {
    geometry, mathtools, array, bm, esint, keytheorems, libertine, inconsolata,
    siunitx, fixdif, derivative, physics2, tasks, anyfontsize
  }
\RequirePackage[inline]{enumitem}
\RequirePackage{fontawesome5, fancyqr, circuitikz}
\hook_gput_code:nnn { cmd/faIcon/before } { . }
  {
    \group_begin:
    \fontsize { \fp_eval:n { .9 * \f@size } } { 0 } \selectfont
  }
\hook_gput_code:nnn { cmd/faIcon/after } { . } { \group_end: }
\geometry{margin = 1in}

\RequirePackage[svgnames]{xcolor}
\newkeytheorem{note}
  [ style = remark, numbered = false,
    name = \color{DarkRed}\sffamily \faIcon{info-circle}~Tip. ]
\newkeytheorem{problem}
  [ style = definition,
    name = \color{DarkGreen}\sffamily \faIcon{pen-square}~Problem ]
\newkeytheorem{solution}
  [ style = definition, numbered = false,
    name = \color{MidnightBlue}\sffamily \faIcon{check-square}~Solution ]
\usephysicsmodule{ab, op.legacy}
\usephysicsmodule{bm-um.legacy}
\settasks { label = \sffamily\Alph*., item-indent = 2.5em, label-offset = .5pt }
\newcommand \true {\addtocounter{task}{1}\faIcon{check-circle}}
\def \@result#1 {\underline{~#1~}}
\def \s@result {\textsf{【\qquad 】}}
\DeclareRobustCommand \result {\@ifstar\s@result\@result}
\keys_define:nn { textfig }
  {
    ratio      .tl_set:N    = \l__textfig_ratio_tl,
      ratio    .initial:n   = 7 \c_colon_str 2,
    position   .tl_set:N    = \l__textfig_position_tl,
      position .initial:n   = c,
    unknown    .code:n      = \tl_if_novalue:nF {#1}
      { \tl_set_eq:NN \l__textfig_position_tl \l_keys_key_tl }
  }
\NewDocumentEnvironment{textfig}{ O{} }
  {
    \group_begin:
    \keys_set:nn { textfig } {#1}
    \seq_set_split:Nee \l__whats_textfig_ratio_seq { \c_colon_str }
      { \l__textfig_ratio_tl }
    \noindent
    \minipage [ \l__textfig_position_tl ]
      {
        \fp_eval:n
          {
            .96*\seq_item:Nn \l__whats_textfig_ratio_seq {1}/
              (
                \seq_item:Nn \l__whats_textfig_ratio_seq {1} + 
                \seq_item:Nn \l__whats_textfig_ratio_seq {2}
              )
          }\dimexpr \linewidth
      }
  }
  {
    \endminipage
    \vspace{4.5pt plus 1.5pt minus 1.5pt} \par
    \group_end:
  }
\NewDocumentCommand \textfigsplit {}
  {
    \endminipage
    \hspace{.02\linewidth}
    \minipage [ \l__textfig_position_tl ]
      {
        \fp_eval:n
          {
            .96*\seq_item:Nn \l__whats_textfig_ratio_seq {2}/
              (
                \seq_item:Nn \l__whats_textfig_ratio_seq {1} + 
                \seq_item:Nn \l__whats_textfig_ratio_seq {2}
              )
          }\dimexpr \linewidth
      }
  }

% \hook_gput_code:nnn { cmd/maketitle/before } { . } { \vskip -\baselineskip }
% \hook_gput_code:nnn { cmd/maketitle/after } { . }
%   { \__weyoung_maketitle_aftercode_aux: }
\cs_new_protected:Nn \__weyoung_maketitle_aftercode_aux:
  {
    \vspace*{-2\baselineskip}
    \begin{center}
      \begin{minipage}{.66\textwidth}
        \centering
        \begin{tasks}
          [ label-width = 5\ccwd, item-indent = 2cm,
            after-item-skip = 0pt ](2)
          \task[\textbf{课程信息}：] 理学院/A0715022
          \task[\textbf{授课教师}：] 陈许敏
          \task[\textbf{解析提供}：] \textsf{Skyrmion}
          \task[{\makebox[5em][s]{\textbf{\textsf\TeX~排版}：}}] \textsf{Skyrmion}
        \end{tasks}
      \end{minipage}
      \begin{minipage}{.15\textwidth}
        \centering
        \begin{tikzpicture}
          \node [ inner~sep = 0 ] (wechat) {
            \fancyqr
              [ size = 4em, image~padding = 1,
                image = \raisebox{4em/33}{\textcolor{DarkGreen}{\faIcon{weixin}}},
                r~color = DarkGreen, version = 4
              ]{http://weixin.qq.com/r/sRJiemTEGjllrSAZ90ej}
            };
          \node [ inner~sep = .5ex, align = center, right, font = \scriptsize ]
            at (wechat.east) {未\\央\\学\\社};
        \end{tikzpicture}
      \end{minipage}
      \begin{minipage}{.15\textwidth}
        \centering
        \begin{tikzpicture}
          \node [ inner~sep = 0 ] (qq) {
            \fancyqr
              [ size = 4em, image~x~padding = 2, image~y~padding = 1,
                image = \raisebox{4em/33}{\textcolor{MidnightBlue}{\faIcon{qq}}},
                l~color = MidnightBlue, version = 4
              ]{https://qm.qq.com/q/PH7CsQWvm2}
            };
          \node [ inner~sep = .5ex, align = center, left, font = \scriptsize ]
            at (qq.west) {物\\理\\营};
        \end{tikzpicture}
      \end{minipage}
    \end{center}
  }

\newcommand \iu {{\mathup i\mkern1mu}}
\newcommand \eu {\mathup e}
\AtBeginDocument{%
  \let \Re \relax \let \Im \relax
  \DeclareMathOperator{\Re}{Re}
  \DeclareMathOperator{\Im}{Im}
}

\makeatother \ExplSyntaxOff
\usetikzlibrary{ fadings, patterns, arrows, decorations.markings }
\tikzset
  {
    > = stealth',
    every picture/.append style =
      { line join = round, line cap = round, thick }
  }

\setCJKmainfont{LXGW WenKai}[AutoFakeSlant]
\setCJKsansfont{LXGW Marker Gothic}[AutoFakeBold, AutoFakeSlant, Scale = .95]
\setCJKmonofont{LXGW WenKai Mono}[AutoFakeBold, AutoFakeSlant]
\ctexset{section/format+ = \raggedright}
\usepackage[warnings-off={mathtools-colon,mathtools-overbracket}]{unicode-math}
\setmathfont[AutoFakeBold]{STIX Two Math}

\makeatletter \def \@result #1{\underline{ \quad }} \makeatother
\DeclareMathOperator{\DFT}{DFT}

\title{\sffamily\bfseries 24Fall \textsc{Hdu} 「数字信号处理」期末试题}
\author{} \date{}

\begin{document}

\maketitle

\section{选择题（每题2分，共20分）}

\begin{problem}
  已知系统单位冲击响应函数 $h(n) = 2^nu(-n)$，则该系统 \result{}（是 / 不是） 稳定系统.
\end{problem}

\begin{problem}
  已知周期序列为 $\cos(n\pi/5)$，则求该周期序列的 DFS 系数 $X(k)$ 为 \result{}.
\end{problem}

\begin{problem}
  已知一个滤波器的幅频响应函数为
  $|H(e^{j\omega})| =
  \begin{cases}
    0, & |\omega| \leq \omega_c\\
    1, & \omega_c < |\omega| \leq \pi
  \end{cases}$，
  则新的系统单位冲击响应为 $(-1)^nh(n)$，
  则该新系统是 \result{} 通（低通 / 高通 / 带通 / 带阻）滤波器.
\end{problem}

\begin{problem}
  已知时域信号 $R_4(n)$，求该信号的 $4$ 点 $\DFT$ 为 \result{}.
\end{problem}

\begin{problem}
  已知线性相位 FIR 滤波器单位冲击响应函数 $h(n) = \delta(n) - \delta(n - 2)$，
  是第 \result{} 线性相位 FIR 数字滤波器，系统的群时延为 \result{}.
\end{problem}

\begin{problem}
  假设 $x(n) = \{1,2,3,4,n = 0,1,2,3\}$，
  $r(n) = \sum_{r = -\infty}^{+\infty} x(n + 3r)$，$r(n)$ 的周期为 \result{}，
  且 $r(n)$ 在一个周期的值为 \result{}.
\end{problem}

\begin{problem}
  系统差分方程 $2y(n) + 0.6y(n - 1) - 0.4y(n - 2) = 3x(n) + 4.2x(n - 1)$，
  则其滤波器的结构中至少需要延时器的个数为 \result{}.
\end{problem}

\begin{problem}
  系统单位冲击响应 $h(n)$ 为实因果序列，系统函数 $H(z)$ 仅有一个零点且在原点有两个极点，
  其中一个极点为 $z = e^{j\pi/2}$，$H(1) = 0.5$，求 $H(z) =$ \result{}.
\end{problem}

\section{解答题（共80分）}

\begin{problem}[12 pt]
  假设一个因果系统的差分方程为
  \[y(n) = x(n) - 0.9y(n - 1)\]
  \begin{enumerate}
    \item 求出系统函数 $H(z)$，单位冲击响应函数 $h(n)$ 及系统的稳定性.
    \item 写出系统的频率响应函数，分析频率在 $0 \sim \pi$，$\pi \sim 2\pi$
    内的幅频响应曲线变换规律，并说明滤波器特性（低通 / 高通 / 带通 / 带阻滤波器）
    \item 当输入信号为 $x(n) = \cos(\omega_0n)$ 时，求稳态输出响应 $y(n)$.
  \end{enumerate}
\end{problem}

\begin{problem}[10 pt]
  以 $\qty{20}\kHz$ 的采样频率对最高频率为 $\qty{10}\kHz$ 的带限信号 $x(t)$ 进行采样，
  然后计算 $x(n)$ 的 $N = 1000$ 个采样点的 $\DFT$，即
  \[X(k) = \DFT[x(n)] = \sum_{n = 0}^{N - 1} x(n) W_N^{nk}\]
  \begin{enumerate}
    \item 求 $k = 150$ 对应的模拟频率是多少，$k = 800$ 呢？
    \item 求频谱采样点间的频率间隔是多少？
    \item 假如对 $x(n)$ 进行后边补零然后进行 $\DFT$ 运算是为了解决什么问题，
    序列的 $N$ 点 $\DFT$ 和 DTFT 是什么关系.
  \end{enumerate}
\end{problem}

\begin{problem}[12 pt]
  已知序列 $x(n) = \{3,1,-1,2;\ n = 0,1,2,3,4\}$，
  $y(n) = \{1,2,-1;\ n = 0,1,2\}$.
  \begin{enumerate}
    \item 计算 $x(n)$ 与 $y(n)$ 的线性卷积.
    \item 计算 $x(n)$ 与 $y(n)$ 的 $5$ 点的圆周卷积.
    \item 计算 $x(n)$ 与 $y(n)$ 的 $7$ 点的圆周卷积.
    \item 圆周卷积和线性卷积的关系.
  \end{enumerate}
\end{problem}

\begin{problem}[10 pt]
  已知 FIR 滤波器的单位冲击响应函数
  \[h(n) = \ab(\frac12)^n [u(n) - u(n - 5)]\]
  \begin{enumerate}
    \item 求该滤波器的直接型结构图.
    \item 求改波器是否为线性相位.
  \end{enumerate}
\end{problem}

\begin{problem}[14 pt]
  $\DFT$ 的运算流图
  \begin{center}
    \begin{tikzpicture}
      \node [ inner sep = 1.5, fill, circle ] (a) at (-2.5,1) {}
       node [ left = 1.5 ] at (a) {$X_1(k)$};
      \node [ inner sep = 1.5pt, fill, circle ] (b) at (-2.5,-1) {}
       node [ left = 1.5 ] at (b) {$X_2(k)$};
      \node [ inner sep = 1.5, fill, circle ] (c) at (2.5,1) {}
       node [ right = 1.5 ] at (c) {$X_1(k) + W_N^k X_2(k)$};
      \node [ inner sep = 1.5, fill, circle ] (d) at (2.5,-1) {}
       node [ right = 1.5 ] at (d) {$X_1(k) - W_N^k X_2(k)$}
       node [ below left ] at (d) {$-1$};
      \node [ inner sep = 1.5, fill, circle ] (e) at (-.5,1) {};
      \node [ inner sep = 1.5, fill, circle ] (f) at (-.5,-1) {}
       node [ above left ] at (f) {$W_N^k$};
      \draw [ -> ] (a) -- (c);
      \draw [ -> ] (b) -- (f);
      \draw [ -> ] (f) -- (d);
      \draw [ -> ] (e) -- (d);
      \draw [ -> ] (f) -- (c);
    \end{tikzpicture}
  \end{center}
  \begin{enumerate}
    \item 试画出 $4$ 点的基 $2$ 时间抽取法 FFT 的基本蝶形运算流图，
    并且利用所画流程图计算每一级的输出结果，输入信号为 $x(n) = \{x(0),x(1),0,0\}$.
    \item 试说明将时域两个序列 $\{x(0),x(1)\}$ 和 $\{x(0),x(1),0,0\}$
    分别作为 FFT 输入，两个 FFT 输出在哪些值是相等的.
    \item 不计 $\pm j$ 及 $\pm 1$ 的乘法次数，完成 $8$ 点的基 $2$ 按时域抽取的 FFT 所需的复数加法和复数乘法次数.
  \end{enumerate}
\end{problem}

\begin{problem}[12 pt]
  已知设计一个一阶数字低通滤波器，其通带截止频率为 $\qty{200}\Hz$，通带最大衰减为 $\qty3\dB$，采样频率为 $\qty{1000}\Hz$，其巴特沃斯滤波器系统函数为
  \[H_a(s) = \frac{1}{1 + s/\Omega_p}\]
  其中 $\Omega_p$ 通带截止频率单位为 $\unit{\radian/\s}$，求
  \begin{enumerate}
    \item 利用冲邀响应不交法设计 IIR 滤波器 $H(z)$.
    \item 试说明冲激响应不变法设计滤波器的优缺点.
    \item 请问设计的数字滤波器能否是因果稳定的，如果能满足因果稳定需要满是什么条件？
  \end{enumerate}
\end{problem}

\begin{problem}[10 pt]
  用频率采样法设计一个低通 FIR 线性相位滤波器，其截止频率是抽样频率的 $1/20$，采样点数为 $N = 21$，边沿过渡带上设计一个过渡点，且该过渡点的幅度值为 $|H(k)| = 0.25$，试求各采样点的 $H(k)$.
\end{problem}

\end{document}